% This poster was created using the baposter class,
% it was presented at the MoRePaS 2018 international
% conference, and can be viewed online here:
% https://doi.org/10.14293/P2199-8442.1.SOP-MATH.JHXASX.v1

\documentclass[portrait,a0paper]{baposter}

% CUstom packages
\usepackage[utf8]{inputenc}
\usepackage[english]{babel}
\usepackage{lipsum}
\usepackage{multicol}
\usepackage{graphbox}
\usepackage{tikz}
\usepackage{pgfplots}
\pgfplotsset{compat=1.9}
\usepackage[font=small,labelfont=bf,justification=centering]{caption}
\usepackage{natbib}
\usepackage{amsfonts}
\usepackage{amsmath}
\usepackage{stackengine}
\usepackage{siunitx}
\usepackage{booktabs}
\usepackage{enumitem}
\usepackage{tabularx}
\usepackage[none]{hyphenat}
\usepackage{mathtools}
\mathtoolsset{showonlyrefs}

% CUstom colors
\definecolor{mDarkBrown}{HTML}{604c38}
\definecolor{mDarkTeal}{HTML}{23373b}
\definecolor{mLightBrown_old}{HTML}{EB811B}
\definecolor{mLightGreen}{HTML}{14B03D}
\definecolor{mDarkBlue}{HTML}{117777}
\definecolor{boxColor}{HTML}{CCDDAA}
\colorlet{boxColor}{black!2}

\colorlet{mLightBrown}{mLightBrown_old!90!black}

% CUstom font
\usepackage[sfdefault]{FiraSans}
\usepackage{FiraMono}
\makeatletter
\newcommand{\globalcolor}[1]{%
	\color{#1}\global\let\default@color\current@color
}
\makeatother
\AtBeginDocument{\globalcolor{mDarkTeal!80!black}}

% CUstom commands
\newcommand{\alert}[1]{\textcolor{mLightBrown!80}{#1}}
\newcommand{\var}[1]{\textcolor{mDarkBlue!75!black}{#1}}
\usetikzlibrary{plotmarks,trees,arrows,shapes,positioning}
\tikzset{>=latex}
\tikzstyle{active}=[circle,fill=mDarkBlue,minimum size=.25cm]
\captionsetup{labelformat=empty}
\newcommand{\T}{^\mathsf T}
\setlength{\bibsep}{0pt plus 0.5ex}
\stackMath
\sisetup{scientific-notation = fixed,fixed-exponent=0}
\newtagform{eqColor}{\color{mDarkBlue}(}{)}
\usetagform{eqColor}
\sisetup{group-separator = {}}
\newcolumntype{Y}{>{\centering\arraybackslash}X}

% SAve space in lists. Use this after the opening of the list
\newcommand{\compresslist}{%
\setlength{\itemsep}{2pt}%
\setlength{\parskip}{0pt}%
\setlength{\parsep}{0pt}%
}
\setlist[itemize]{leftmargin=12pt}

% PGFPLOTS OPTIONS
\pgfplotsset{every axis/.append style={
		label style={font=\Huge},
		tick label style={font=\small}  
	}}

\begin{document}
\begin{poster}%
%% Poster options
{
	grid=false,
	columns=3,
	colspacing=1em,
	bgColorOne=white,
	headerColorOne=mLightBrown!0,
	headerFontColor=mLightBrown!80,
	borderColor=mLightBrown!50,
	textborder=rectangle,
	eyecatcher=true,
	headerborder=none,
	headerheight=0.1\textheight,
	boxheaderheight=1.75em,
	headershape=smallrounded,
	headershade=plain,
	headerfont=\large\bfseries, %Sans Serif
	boxColorOne=boxColor,
	background=plain,
	linewidth=1pt
}
	% Eye Catcher
{

}
	% Title
{\spaceskip 0.4 em\relax Statistical Learning in Tree-Based\linebreak Tensor Format\vspace{0.5em}}
	% Authors
{\hfill {Erwan Grelier}\textsuperscript{\dag} \hfill {Mathilde Chevreuil} \hfill {Anthony Nouy} \hfill \vphantom{A} \vspace{-5mm}}
{
	\begin{tabular}{cccccc}
		\\[7mm]
		\multicolumn{3}{c}{\includegraphics[align=c,height=1.3cm]{./doc/logo_centrale.eps}} &
		\multicolumn{3}{c}{\includegraphics[align=c,height=1.3cm]{./doc/logo_univ.eps}} \\[.6cm]
		\multicolumn{2}{c}{\includegraphics[align=c,height=1.3cm]{./doc/logo_lmjl.pdf}} &
		\multicolumn{2}{c}{\includegraphics[align=c,height=1.6cm]{./doc/logo_GeM.eps}} &
		\multicolumn{2}{c}{\includegraphics[align=c,height=1.3cm]{./doc/logo_CNRS.eps}}
	\end{tabular}
}

%% Boxes
\headerbox{}{name=affiliations,column=0,row=0,span=3,boxheaderheight=0em,boxColorOne=white,headerColorOne=white,textborder=none,headerborder=none}{
	\vspace{-.4cm}
	\begin{center}
		\begin{tabularx}{\textwidth}{YYYY}
			\hspace{4.5mm}Centrale Nantes&
			Université de Nantes&
			\hspace{-6.9mm}Centrale Nantes&
			 \\
			\hspace{4.5mm}GeM (UMR CNRS 6183)&
			GeM (UMR CNRS 6183)&
			\hspace{-6.9mm}LMJL (UMR CNRS 6629)&
			\mbox{\textsuperscript{\dag}\textbf{\texttt{erwan.grelier@ec-nantes.fr}}}
		\end{tabularx}
	\end{center}
}

\headerbox{References}{name=references,column=0,span=2,above=bottom,textborder=none,headerColorOne=white,headerborder=none,headerFontColor=mDarkTeal!80!black}{
	\nocite{Hackbusch2012}
	\nocite{Grasedyck2010}
	\nocite{nouy:hal-01262422}
	\nocite{Chevreuil1}
	\renewcommand{\section}[2]{}
	\bibliographystyle{abbrv}
	\bibliography{bibliography}
}

\headerbox{\begin{tabular}{c}
Statistical learning for \\uncertainty quantification
\end{tabular}}{name=UQ,column=0,below=affiliations,span=1,boxheaderheight=3.3em}{
	We denote by
	\vspace{-2mm}
	\begin{itemize}\compresslist
		\item $X = (X_1,\ldots,X_d)$ a set of \alert{$d$ random variables},
		\item $Y = f(X)$ the random output of a \alert{\mbox{model $f$}, costly to evaluate}.
	\end{itemize}
	The \alert{uncertainty quantification} requires many realizations of $Y$, hence \alert{many evaluations of $f$}. To reduce their cost we use an \alert{approximation $v$ of $f$}, by solving
	\begin{equation}
		\label{eq: problem}
		\min\limits_{\var{v} \in \mathcal V} \frac 1 N \sum_{k=1}^N \left|y^k - \var{v}(x^k)\right|^2
	\end{equation}
	with $\mathcal V$ an approximation set and $(x^k,y^k)$ a \alert{realization of $(X,Y)$}.
}

\headerbox{Tree-based (TB) tensor format}{name=TB,column=0,below=UQ,above=references}{
	A function \alert{$v \in \mathcal T^T_r$} in \alert{tree-based \mbox{tensor format}} associated to a \alert{dimension partition tree} \mbox{$T\subset 2^{\{ 1,\ldots,d \}}$}
	\vspace{-10mm}
	\begin{center}
		\footnotesize
\begin{tikzpicture}[scale=1.2,level distance = 4mm,baseline=(current bounding box.south),label distance=1mm]
\tikzstyle{level 1}=[sibling distance=15mm]
\tikzstyle{level 2}=[sibling distance=8mm]

\node [active,label={[label distance=0mm]above:{$\var{f_{1,2,3,4}}$}}]  {}
child {node [active,label={[label distance=0mm]above left:{$\var{f_{1,2}}$}}]  {}
	child { node [active,label=below:{$\var{f_{1}}$}] {}} 
	child { node [active,label={below:$\var{f_{2}}$}] {}}
	}
child {node [active,label={[label distance=0mm]above right:{$\var{f_{3,4}}$}}] {}
	child { node [active,label=below:{$\var{f_{3}}$}] {}}
	child { node [active,label=below:{$\var{f_{4}}$}] {}}  
	}
;
\end{tikzpicture}
	\end{center}
	\vspace{-4mm}
	admits a representation as a \alert{composition of multilinear vector-valued functions} (tensors):
	\begin{alignat*}{2}
		v(x) &= v(x_1,x_2,x_3,x_4) \\ 
		&= \var{f_{1,2,3,4}}(
			\underbrace{\var{g_{1,2}}(x_1,x_2)}_{\var{f_{1,2}}(\var{f_1}(x_1),\var{f_2}(x_2))},\underbrace{\var{g_{3,4}}(x_3,x_4)}_{\var{f_{3,4}}(\var{f_3}(x_3),\var{f_4}(x_4))})
	\end{alignat*}
	where for the leaves of $T$, $1 \leq \nu \leq d$, 
	\begin{equation*}
	\var{f_\nu} : \mathcal X^\nu \rightarrow \mathbb R^{r_\nu},
	\end{equation*}
	and for a node $\alpha$ with $s$ children $S(\alpha)$,
	\begin{equation*}
	\var{f_\alpha} : \bigtimes_{\mathclap{\beta \in S(\alpha)}} \mathbb R^{r_{\beta}} \rightarrow \mathbb R^{r_\alpha}
	\end{equation*}
	is a multilinear function identified with a tensor in $\mathbb R^{r_\alpha \times r_{\beta_1} \times \cdots \times r_{\beta_s}}$.
	\bigbreak
	Properties of the set $\mathcal T^T_r$:
	\vspace{-2mm}
	\begin{itemize}\compresslist
		\item \alert{closed}, best approximation problems are well posed,
		\item \alert{storage complexity} of $v$ scaling linearly \mbox{with $d$},
		\item \alert{multilinear parametrization} of $v$,
		\item existence of a \alert{higher-order singular value decomposition} (HOSVD) of $v$.
	\end{itemize}
}

\headerbox{Statistical learning in TB tensor format}{name=SLTB,column=1,below=affiliations}{
	Using the multilinear parametrization of \mbox{$v \in \mathcal T^T_r$}, problem \eqref{eq: problem} is solved using an \alert{alternating minimization algorithm}, leading to \alert{linear problems}:
	\begin{equation*}
		\min\limits_{\var{a_\alpha} \in \mathbb R^{m_\alpha}} \frac 1 N \sum\limits_{k=1}^N \left| y^k - \Psi_\alpha(x^k)\T \var{a_\alpha} \right|^2,\; \forall \alpha \in T
	\end{equation*}
	with $\Psi_\alpha(x)$ such that $v(x) = \Psi_\alpha(x)\T \var{a_\alpha}$. For the leaves of $T$, a \alert{sparse approximation} is computed by using a \alert{working-set} strategy.
	\bigbreak
	\textbf{Iterative adaptation of $r$} 
	
	Starting from $r^0 = (0,\ldots,0)$, at iteration $n$, given an approximation $v^n \in \mathcal T^T_{r^n}$:
	\begin{itemize}\compresslist
		\item \alert{increase the ranks} $\left( r^n_\beta \right)_{\beta \in T_n}$ associated with a \alert{subset of nodes $T_n$}$ \subset T$,		
		\item $T_n$ is chosen as
		\begin{equation*}
			T_n = \left\{ \alpha \in T: \sigma_\alpha \ge \theta \max\limits_{\beta \in T} \sigma_\beta \right\},
		\end{equation*}
		with $\theta \in [0,1]$ and $\sigma_\alpha$ the smallest \alert{$\alpha$-singular value} for the node $\alpha$, obtained by computing the SVD of an $\alpha$-matricization of $v^n$.
	\end{itemize}
}

\headerbox{First numerical experiment}{name=results,column=1,below=SLTB,above=references}{
	Approximation of 
	\begin{equation*}
		f(X) = \left( 1 + \sum_{i=1}^d a_i X_i \right)^{-2},
	\end{equation*}
	\begin{itemize}\compresslist
		\item $X = (X_1,\ldots,X_{10})$, $X_i \sim \mathcal U(0,1)$,
		\item $0 \leq a_i \leq 1$, $i = 1,\ldots,d$,
		\item \alert{linear tree} (tensor train Tucker format),
		\item approximation spaces at the leaves: polynomial spaces of \alert{maximal degree \num{20}},
		\item \alert{\num{10000} training samples}.
	\end{itemize}
	\vspace{2mm}
	\hspace{-5.75mm}
	\includegraphics[scale=1]{./doc/cornerPeakHTT.pdf}
}

\headerbox{TB learning and changes of variables}{name=TBCOV,column=2,below=affiliations}{
	\newcounter{treestep}
	\setcounter{treestep}{0}
	An approximation is searched in the form
	\vspace{-1mm}
	\begin{equation}
		v(x) = \var{h}(\var{g}(x)) = \var{h}(\var{g_1}(x),\ldots,\var{g_m}(x)),
	\end{equation}
	with $\var{h} \in \mathcal T^T_r$ and $\var{g_i}: \mathbb R^d \rightarrow \mathbb R$, $i = 1,\ldots,m$.
	\bigbreak
	\textbf{Outline of the algorithm}

 	Construction of approximations $v_m$, at each iteration \alert{introducing a new variable} \mbox{$\var{z_m} =\var{g_m}(x)$}, and alternatively optimizing on
 	\begin{itemize}\compresslist
 		\item $\var{h} \in \mathcal T^{T_m}_{r_m}$, with fixed $\var g$, using the \alert{statistical learning algorithm in TB tensor format},
 		\item $\var{g}$, with fixed $\var h$, using a \alert{Gauss-Newton algorithm},
 	\end{itemize}
 	\vspace{-2mm}
 	until convergence.

 	\vspace{-1.65cm}
 	\hspace{-.75cm}
 	\hfill
	\begin{minipage}{.2\textwidth}
		\begin{center}
			\footnotesize
\begin{tikzpicture}[scale=1,sibling distance=9mm,level distance = 8mm,baseline=(current bounding box.south),label distance=1mm]

\draw[draw=none] (0,0) -- ++ (-4mm*\thetreestep,-8mm*\thetreestep) -- ++ (-4mm,-4mm);
\draw[draw=none] (0,0) -- ++ (8mm-4mm*\thetreestep,16mm-8mm*\thetreestep) -- ++ (8mm,8mm);

\node [active,label=above:{$1$}]  {}
child {node [active,label=below:{$\stackon{\alert{\{z_1\}}}{2}$}]  {}}
child {node [active,label=below:{$\stackon{\alert{\{z_2\}}}{3}$}] {}}
;
\end{tikzpicture}
		\end{center}
	\end{minipage}
	\hspace{.15cm}
	\begin{minipage}{.2\textwidth}
		\begin{center}
			\footnotesize
\begin{tikzpicture}[scale=1,sibling distance=9mm,level distance = 8mm,baseline=(current bounding box.south),label distance=1mm]

\addtocounter{treestep}{1}
\draw[draw=none] (0,0) -- ++ (-4mm*\thetreestep,-8mm*\thetreestep) -- ++ (-4mm,-4mm);
\draw[draw=none] (0,0) -- ++ (8mm-4mm*\thetreestep,16mm-8mm*\thetreestep) -- ++ (8mm,8mm);

\node [active,label=above:{$1$}]  {}
child {node [active,label=above left:{$2$}]  {}
	child { node [active,label=below:{$\stackon[1pt]{\alert{\{z_1\}}}{4}$}] {}}
	child { node [active,label=below:{$\stackon[1pt]{\alert{\{z_2\}}}{5}$}] {}}
}
child {node [active,label=below:{$\stackon[1pt]{\alert{\{z_3\}}}{3}$}] {}}
;
\end{tikzpicture}
		\end{center}
	\end{minipage}
	\hspace{.5cm}
	\begin{minipage}{.2\textwidth}
		\begin{center}
			\include{./doc/sequentialTTT3}
		\end{center}
	\end{minipage}
	\hfill
	\vspace{-.25cm}
}

\headerbox{Second numerical experiment}{name=results_2,column=2,below=TBCOV}{
	Approximation of 
	\begin{align}
		f(X) = &\sin(\mathrm w_2\T X + \mathrm w_3\T X) \cos(\mathrm w_1\T X + \mathrm w_4\T X) + \notag \\
			&\cos(\mathrm w_1\T X + \mathrm w_3\T X),
	\end{align}
	\begin{itemize}\compresslist
			\vspace{-5mm}
		\item $d = 20$,
		\item $X = (X_1,\ldots,X_d)$, 
		\item $X_i \sim \mathcal U(-1,1)$, $i = 1,\ldots,d$,
		\item \alert{$m \leq 4$},
		\item approximation spaces at the leaves: polynomial spaces of \alert{maximal degree \num{10}},
		\item approximation space for the $\var{g_j}$, ${j = 1,\ldots,4}$: \alert{linear functions},
		\item \alert{$N = \num{1000}$ training samples}.
	\end{itemize}
	\begin{center}
		\begin{tabular}{ccc} \toprule
			$v$ & \begin{tabular}{@{}c}Confidence interval\\ for the error\end{tabular} & \begin{tabular}{@{}c}Storage\\complexity\end{tabular} \\ \midrule
			$\var{h}$ & $\left[\num{2e-5},\num{1.1} \right] \cdot 10^{-1}$ & \num{435} \\
			$\var{h}\circ \var{g}$ &  
			$\left[\num{1.6},\num{7.9} \right] \cdot 10^{-3}$ & \num{342} \\\bottomrule
		\end{tabular}
	\end{center}
	\medbreak
	$N$ being small compared to $d$, introducing changes of variables \alert{reduces the mean and variance of the error, and the storage complexity}.
}

\headerbox{Conclusions}{name=conclusion,column=2,span=1,below=results_2}{
	The statistical learning in tree-based tensor format:
	\vspace{-.1cm}
	\begin{itemize}\compresslist
		\item uses the \alert{classical machinery of linear approximation},
		\item exploits both \alert{low-rank} and \alert{sparsity},
		\item can be combined with \alert{changes of variables}.
	\end{itemize}
}

\headerbox{Outlooks}{name=outlooks,column=2,span=1,below=conclusion,above=bottom}{
	\begin{itemize}\compresslist
		\item Perform \alert{tree adaptation},
		\item use \alert{other trees} for $\var h \in \mathcal T^T_r$ when $v = \var h \circ \var g$,
		\item \alert{characterize the properties of the set} of functions of the form $\var h\circ \var g$, $\var h \in \mathcal T^T_r$.
	\end{itemize}
}
\end{poster}
\end{document}